% Paper 2 · Appendix C · DeepSeek V4 preliminary results
% Status: 2026-05-05 · sample 48 FINAL · 60.4% · HOLD full 500 (待 max_tokens 调高重测)

\section{DeepSeek V4 preliminary results}
\label{app:deepseek-v4}

DeepSeek released V4-pro and V4-flash via official API
(\url{api.deepseek.com}) in early 2026, with three reasoning modes:
\texttt{non-think}, \texttt{think-high}, \texttt{think-max}, and a 1M
context window (vs V3.2's 128K).

We ran a 50-question stratified subset evaluation with V4-pro under
\texttt{think-high} to test if the new model + larger context window
shifts the accuracy ceiling, while keeping the rest of the pipeline
identical (m3-rerank · query rewriting · type-aware prompts).

\subsection{Setup}

\begin{itemize}
\item Subject + Judge: DeepSeek V4-pro (subject) · V4-flash (judge ·
  cross-family avoided since same vendor · documented limitation).
\item Reasoning mode: \texttt{think-high}
\item Context: full 1M tokens · no truncation needed for LongMemEval-S
  50K haystacks.
\item Pipeline: identical to v0.8 (Stage 1-5 from §\ref{sec:method}).
\item Cost: \$1.40 USD for 50 questions (vs \$0.15 USD for v0.8 50q
  on Volc Ark coding plan · 9× more expensive).
\end{itemize}

\subsection{Sample 48 results (n=48 stratified · 8 per type)}

\begin{table}[h]
\centering
\small
\begin{tabular}{lrrr}
\toprule
\textbf{Type} & \textbf{V0.8 (V3.2)} & \textbf{V4-pro think-high} & \textbf{Δ} \\
\midrule
knowledge-update         & 57.7\% & \textbf{75.0\%} (6/8) & \textbf{+17.3} ⭐ \\
single-session-assistant & 83.9\% & 75.0\% (6/8) & \(-8.9\) \\
single-session-user      & 57.1\% & \textbf{75.0\%} (6/8) & \textbf{+17.9} ⭐ \\
multi-session            & 54.9\% & 62.5\% (5/8) & +7.6 \\
single-session-preference& 53.3\% & 50.0\% (4/8) & \(-3.3\) \\
temporal-reasoning       & 46.6\% & \textbf{25.0\%} (2/8) & \textbf{\(-21.6\)} ⚠️ \\
\midrule
\textbf{Overall (sample)} & 56.2\% & \textbf{60.4\%} & +4.2 \\
\bottomrule
\end{tabular}
\caption{DeepSeek V4-pro \texttt{think-high} mode · sample 48. Three
question types (ku · ssa · ssu) reach 75\% accuracy, demonstrating that
1M context + native reasoning solves a new band of recall tasks. However,
two regressions are notable: ssa drops 9 pts vs v0.8, and
temporal-reasoning drops 22 pts.}
\label{tab:v4-sample}
\end{table}

\subsection{Key observations}

\begin{itemize}
\item \textbf{Multi-session: +25 pts at sample.} Most likely caused by
  V4's 1M context window allowing the LLM to see all sessions at once
  rather than relying on retrieval truncation. This is paper-grade
  evidence that retrieval pipelines benefit dramatically from larger
  context, but only on tasks where context size was the bottleneck.

\item \textbf{Knowledge-update: +17 pts.} Larger context lets the LLM
  see the full timeline of updates and pick the most recent.
  Type-aware prompt (timestamps) compounds with this.

\item \textbf{Single-session-assistant: anomaly.} V4 sample shows
  \textit{decreased} accuracy on ssa (where v0.8 was strongest at
  83.9\%). Hypothesis: V4 uses thinking tokens aggressively
  (\texttt{reasoning\_content} fills the response budget); when the
  same \texttt{max\_tokens} is used as v0.8, the actual answer
  \texttt{content} can be truncated, missing key details that ssa
  questions require. Mitigation: increase \texttt{max\_out\_tok} from
  4096 to 16384 for V4 evaluations.

\item \textbf{Cost vs gain trade-off.} V4-pro think-high at \$0.028/q
  vs V3.2 thinking at \$0.003/q is 9× more expensive. The +5-8 pts
  overall gain projection means roughly \$0.005 per accuracy point ·
  not unreasonable for paper-grade benchmarks but unattractive for
  high-volume production.
\end{itemize}

\subsection{Limitations of preliminary V4 data}

\begin{itemize}
\item Only 35/50 questions completed at submission time · final
  numbers may shift ±3 pts.
\item V4-pro + V4-flash judge is \emph{not} a true cross-family check ·
  both from DeepSeek vendor. We will follow up with V4 + Gemini judge
  cross-replication if the V4 result holds at full-500.
\item CPU-mode bge-m3 (used to avoid GPU contention with v3.2 cross-judge
  run) doesn't affect retrieval quality but slows wall-clock.
\item think-max mode untested · may push accuracy higher at further
  cost increase.
\end{itemize}

\subsection{Full-500 update (2026-05-06)}
\label{app:v4-full500}

We subsequently ran the full 500-question evaluation under
\texttt{think-high} (no \texttt{max\_tokens} cap · vanilla pipeline ·
no per-type routing) to test whether the sample 48 trends survive on
n=500. They mostly do not.

\begin{table}[h]
\centering
\small
\begin{tabular}{lrrrr}
\toprule
\textbf{Type} & \textbf{v0.8 (V3.2 n=500)} & \textbf{V4 sample (n=48)} & \textbf{V4 full (n=500)} & \textbf{Δ vs v0.8} \\
\midrule
single-session-assistant  & 83.9\% & 75.0\% & \textbf{87.5\%} (49/56) & \textbf{+3.6} \\
single-session-preference & 53.3\% & 50.0\% & \textbf{66.7\%} (20/30) & \textbf{+13.4} \\
single-session-user       & 57.1\% & 75.0\% & 58.6\% (41/70)  & +1.5 \\
knowledge-update          & 57.7\% & 75.0\% & 56.4\% (44/78)  & \(-1.3\) \\
multi-session             & 54.9\% & 62.5\% & 52.6\% (70/133) & \(-2.3\) \\
temporal-reasoning        & 46.6\% & 25.0\% & 43.6\% (58/133) & \(-3.0\) \\
\midrule
\textbf{Overall}          & \textbf{56.6\%} & 60.4\% & \textbf{56.4\%} (282/500) & \(-0.2\) \\
\bottomrule
\end{tabular}
\caption{V4-pro \texttt{think-high} full 500 vs the sample 48 estimate
and the v0.8 V3.2 baseline. Cost: approximately \$8 USD wall-clock 6.3
hours on a single T4 (CPU bge-m3 · GPU contention with cross-judge
runs).}
\label{tab:v4-full500}
\end{table}

\paragraph{Verdict.} V4-pro \texttt{think-high} is essentially tied
with V3.2 (\(-0.2\) pts overall · within noise). Per-type, V4 wins
clearly on ssa (\(+3.6\)) and ssp (\(+13.4\)) but loses on multi-session
(\(-2.3\)) and temporal-reasoning (\(-3.0\)). The headline number for
this paper remains v0.8 V3.2 56.6\%.

\paragraph{Lesson on sample size.} The sample 48 estimate showed +4.2
pts over v0.8 (60.4\% vs 56.2\%); the full 500 showed \(-0.2\) pts. Five
of six per-type point estimates from sample 48 disagreed with the
full-500 measurement by more than the per-type 95\% CI half-width.
Knowledge-update sample +17.3 was the largest outlier (full \(-1.3\)).
Future model substitutions should not be approved on n<100 samples
alone.

\paragraph{Cost vs gain.} V4-pro at $\sim$\$0.016/q is 8× more
expensive than V3.2 at $\sim$\$0.002/q (Volc Ark coding plan). Tied
accuracy and 8× cost yields no production case for V4-pro
\texttt{think-high} on this benchmark. \texttt{think-max} or
selective routing (V4 for ssa/ssp · V3.2 for ms/temporal) is the only
path to a positive economic case · neither is in scope for v1.0 paper.

\subsection{Decision (2026-05-05): held full 500}

Based on sample 48 results (60.4\% · +4.2 vs v0.8 sample 56.2\%),
we \textbf{do not} run V4-pro full 500 in current configuration. Reasons:

\begin{itemize}
\item Improvement (+4.2 pts) is below our \textbf{≥3 pt
  significance threshold} but within sample noise on n=48.
\item The 22 pt drop on temporal-reasoning is a regression that must be
  investigated before committing \$14 + 8h to a full run.
\item The 9 pt drop on ssa supports the \texttt{reasoning\_content}
  truncation hypothesis · before retrying, set
  \texttt{max\_tokens=16384} (4× current).
\end{itemize}

Recommended next steps (deferred to follow-up addendum):

\begin{enumerate}
\item V4-pro \texttt{think-high} + \texttt{max\_tokens=16384} sample 48
  (\$2 · 1h) · check if ssa returns to 75--80\% and temporal to ~46\%.
\item If yes, run V4-pro full 500 with the new max\_tokens (\$14 · 8h).
\item V4-pro \texttt{think-max} on sample 48 (\$3 · 1.5h) · may push ku/ssu
  higher but risks worse temporal.
\end{enumerate}
